% !PRINTED 135
\weberpage{135}
%<<CONT

also nach (8) und (12)
\begin{equation}
 |\psi(x)|\,|x-\alpha|^a<|\psi_1(x)|\,|x-\alpha_1|^{a_1}\,|x-\alpha_2|^{a_2}\cdots+|\varphi(x)|,
\tag{16}
\end{equation}
und wenn wir nun $x$ auf der Peripherie von $\varrho$ annehmen, also
\[
 |x-\alpha|=\varrho
\]
setzen, so ist
\[
 |x-\alpha_1|<2\varrho,
 \qquad |x-\alpha_2|<2\varrho,
 \ldots
\]
folglich nach (10) und (16)
\begin{equation}
 \varrho^a A <(2\varrho)^{a_1+a_2+\cdots}B+|\varphi(x)|.
\tag{17}
\end{equation}
Wir wollen nun die oberen Grenzen für die Coëfficienten $\varepsilon_1,\varepsilon_2,
\ldots$ von $\varphi$ so klein annehmen, dass
\[
 |\varphi(x)|<\varrho^aA'
\]
wird, worin $A'$ eine Zahl bedeutet, die kleiner als $A$ ist; dann folgt aus (17)\begingroup\renewcommand{\thefootnote}{E1}\footnote[1]{[Quellenhinweis:] Mit $s=a_1+a_2+\cdots$ ergibt (17) die Schranke $A-A'<2^s\varrho^{s-a}B$. Der Widerspruch betrifft $s>a$, nicht $a>s$: Bei $0<\varrho<1$ gilt dann $2^s\varrho^{s-a}B\le 2^nB\varrho\to0$. Daher müssen in (19), (20) die Zeichen $\ge$ stehen; die Summe liefert weiterhin Gleichheit. Die Verweisung (20), (21) meint (19), (20). Drucktext und Druckformeln bleiben unverändert.}\endgroup
\begin{equation}
 A-A'<2^a\varrho^{-a+a_1+a_2+\cdots}B.
\tag{18}
\end{equation}
Dies aber würde für ein hinreichend kleines $\varrho$ nicht mehr möglich sein, wenn $a$ grösser als die Summe $a_1+a_2+\cdots$ wäre. Es folgt also
\begin{equation}
 a\le a_1+a_2+\cdots .
\tag{19}
\end{equation}
Ebenso lässt sich beweisen, dass
\begin{equation}
 b\le b_1+b_2+\cdots,
 \qquad
 c\le c_1+c_2+\cdots,
 \qquad \ldots
\tag{20}
\end{equation}
sein muss. Da aber die Summen der linken Seiten sowohl als der rechten in den Ungleichungen (20), (21) gleich $n$ sein müssen,
% [sic] Source-original mathematical cross-reference: the printed page says (20), (21), although the displayed inequalities are numbered (19), (20).
so können nur die Gleichheitszeichen bestehen, also:
\[
 a=a_1+a_2+\cdots,
 \quad b=b_1+b_2+\cdots,
 \quad c=c_1+c_2+\cdots,
 \quad \ldots,
\]
wodurch auch der zweite Theil unseres Satzes bewiesen ist.

Wir wollen dem bewiesenen Satze noch folgende, auf den Fall mehrfacher Wurzeln bezügliche Bemerkung beifügen.

Die nothwendige und hinreichende Bedingung dafür, dass $f(x)=0$ überhaupt mehrfache Wurzeln habe, ist die, dass $f(x)$ und $f'(x)$ einen gemeinsamen Factor haben. Wenn man also nach den im ersten Abschnitt entwickelten Principien den Algorithmus
%CONT>>
